\documentclass[a4paper, 14pt]{extarticle}

% Поля
%--------------------------------------
\usepackage{geometry}
\geometry{a4paper,tmargin=2cm,bmargin=2cm,lmargin=3cm,rmargin=1cm}
%--------------------------------------


%Russian-specific packages
%--------------------------------------
\usepackage[T2A]{fontenc}
\usepackage[utf8]{inputenc} 
\usepackage[english, main=russian]{babel}
%--------------------------------------

\usepackage{textcomp}

% Красная строка
%--------------------------------------
\usepackage{indentfirst}               
%--------------------------------------             


%Graphics
%--------------------------------------
\usepackage{graphicx}
\graphicspath{ {./images/} }
\usepackage{wrapfig}
%--------------------------------------

% Полуторный интервал
%--------------------------------------
\linespread{1.3}                    
%--------------------------------------

%Выравнивание и переносы
%--------------------------------------
% Избавляемся от переполнений
\sloppy
% Запрещаем разрыв страницы после первой строки абзаца
\clubpenalty=10000
% Запрещаем разрыв страницы после последней строки абзаца
\widowpenalty=10000
%--------------------------------------

%Списки
\usepackage{enumitem}

%Подписи
\usepackage{caption} 

%Гиперссылки
\usepackage[hidelinks]{hyperref}

\hypersetup {
	unicode=true
}

%Рисунки
%--------------------------------------
\DeclareCaptionLabelSeparator*{emdash}{~--- }

%--------------------------------------

\usepackage{tempora}
\usepackage{amsmath}
\usepackage{color}
\usepackage{ragged2e}
\usepackage{xcolor}
\usepackage{listings}
\usepackage{hyperref}
\usepackage{cleveref}
\definecolor{codegreen}{rgb}{0,0.6,0}
\definecolor{codegray}{rgb}{0.5,0.5,0.5}
\definecolor{codepurple}{rgb}{0.58,0,0.82}

\lstset{
    backgroundcolor=\color{white},
    basicstyle=\ttfamily\small,
    breaklines=true,
    frame=single,
    numbers=none,          % Отключить номера строк по умолчанию
    captionpos=b,          % Подпись снизу
    showstringspaces=false % Не показывать пробелы в строках
}
\lstdefinestyle{python}{
    language=Python,
    keywordstyle=\color{blue}, % Ключевые слова (def, class, import и т.д.)
    commentstyle=\color{green!50!black}, % Комментарии (# ...)
    stringstyle=\color{red}, % Строки ("..." и '...')
    emph={self, True, False, None, int, str, float, bool, list, dict, set, tuple},
    emphstyle=\color{orange}
}

% Оформление заголовков разделов в соответствии с ГОСТ 7.32-2017:
% полужирное начертание, без точки после номера, без автоматического отступа
%--------------------------------------
\usepackage{titlesec}
\titleformat{\section}{\bfseries\large}{\thesection}{1em}{}
\titleformat{\subsection}{\bfseries}{\thesubsection}{1em}{}
\titlespacing*{\section}{0pt}{2ex}{1.5ex}
\titlespacing*{\subsection}{0pt}{1.5ex}{1ex}
%--------------------------------------

\begin{document}

%--------------------------------------
%			ТИТУЛЬНЫЙ ЛИСТ
%--------------------------------------
\begin{titlepage}
\thispagestyle{empty}
\newpage


%Шапка титульного листа
%--------------------------------------
\vspace*{-60pt}
\hspace{-65pt}
\begin{minipage}{0.3\textwidth}
\hspace*{-20pt}\centering
\includegraphics[width=\textwidth]{emblem}
\end{minipage}
\begin{minipage}{0.67\textwidth}\small \textbf{
\vspace*{-0.7ex}
\hspace*{-6pt}\centerline{Министерство науки и высшего образования Российской Федерации}
\vspace*{-0.7ex}
\centerline{Федеральное государственное автономное образовательное учреждение }
\vspace*{-0.7ex}
\centerline{высшего образования}
\vspace*{-0.7ex}
\centerline{<<Московский государственный технический университет}
\vspace*{-0.7ex}
\centerline{имени Н.Э. Баумана}
\vspace*{-0.7ex}
\centerline{(национальный исследовательский университет)>>}
\vspace*{-0.7ex}
\centerline{(МГТУ им. Н.Э. Баумана)}}
\end{minipage}
%--------------------------------------

%Полосы
%--------------------------------------
\vspace{-25pt}
\hspace{-35pt}\rule{\textwidth}{2.3pt}

\vspace*{-20.3pt}
\hspace{-35pt}\rule{\textwidth}{0.4pt}
%--------------------------------------

\vspace{1.5ex}
\hspace{-35pt} \noindent \small ФАКУЛЬТЕТ\hspace{80pt} <<Информатика и системы управления>>

\vspace*{-16pt}
\hspace{47pt}\rule{0.83\textwidth}{0.4pt}

\vspace{0.5ex}
\hspace{-35pt} \noindent \small КАФЕДРА\hspace{50pt} <<Теоретическая информатика и компьютерные технологии>>

\vspace*{-16pt}
\hspace{30pt}\rule{0.866\textwidth}{0.4pt}
  
\vspace{11em}

\begin{center}
\Large {\bf Лабораторная работа № 2 } \\ 
\large {\bf по курсу <<Базы данных>>} \\ 
\end{center}\normalsize

\vspace{8em}


\begin{flushright}
  {Студент группы ИУ9-51Б Бондарев Б.А.\hspace*{15pt} \\
  \vspace{2ex}
  Преподаватель Вишняков И. Э.\hspace*{15pt}}
\end{flushright}

\bigskip

\vfill
 

\begin{center}
\textsl{Москва 2026}
\end{center}
\end{titlepage}

\section{Постановка задачи}
Целью данной лабораторной работы является моделирование данных с
использованием модели семантических объектов.
Для достижения цели были поставлены следующие задачи:
\begin{enumerate}
    \item создать модель семантических объектов для предметной области; 
    \item обосновать выбор кардинальных чисел атрибутов и типов объектов.
\end{enumerate}

\section{Практическая реализация}

\subsection{Модель семантических объектов}
В качестве предметной обалсти было выбрано администрирование Minecraft-сервера. Для построения модели, 
представленой на \hyperref[fig:som-model]{рисунке~\ref*{fig:som-model}}, было выделено 4 семантических объекта:

\begin{enumerate}
    \item \textbf{PLAYER} (игрок). Для каждого игрока система должна хранить:
    \begin{itemize}
        \item идентификатор игрока PlayerUUID (1.1);
        \item игровое имя Username (1.1);
        \item дату и время первого появления на сервере FirstSeenAt (1.1);
        \item дату и время последней активности LastSeenAt (1.1);
        \item роль игрока --- объектный атрибут ROLE (1.1);
        \item статистику посещения игровых миров --- групповой многозначный атрибут PlayerWorld (0.N), в состав которого входят:
        \begin{itemize}
            \item посещённый мир --- объектный атрибут WORLD (1.1);
            \item последняя сохранённая позиция игрока --- групповой атрибут LastPosition (1.1), состоящий из простых атрибутов X, Y и Z (1.1);
            \item дата и время последнего посещения мира LastVisitAt (1.1);
            \item суммарное время игры данного игрока в этом мире Playtime (1.1);
        \end{itemize}
        \item блокировки, полученные игроком, --- объектный атрибут BAN в роли Received (0.N);
        \item блокировки, наложенные игроком, --- объектный атрибут BAN в роли Issued (0.N).
    \end{itemize}
    Атрибут PlayerUUID --- уникальный идентификатор объекта, атрибут Username --- неуникальный идентификатор, поскольку игровое имя может быть изменено игроком и затем занято другим игроком. Атрибут WORLD является идентификатором внутри группы PlayerWorld, то есть один и тот же мир встречается в статистике игрока не более одного раза. Объект PLAYER является гибридным: групповой атрибут PlayerWorld содержит в себе объектный атрибут WORLD. В отличие от модели <<сущность-связь>>, где эта информация хранилась в сущности <<Игровой мир>>, статистика посещения оформлена групповым атрибутом, а не отдельным объектом, так как запись о времени игры и позиции игрока не может существовать независимо от игрока и мира и полностью определяется парой <<игрок --- мир>>.

    \item \textbf{WORLD} (мир). Для каждого игрового мира система должна хранить:
    \begin{itemize}
        \item уникальное название мира WorldName (1.1);
        \item seed генерации мира Seed (1.1);
        \item режим игры GameMode (1.1);
        \item размер границы мира BorderSize (1.1);
        \item измерение (тип) мира Dimension (1.1);
        \item игроков, посещавших мир, --- объектный атрибут PLAYER (0.N).
    \end{itemize}
    Атрибут WorldName --- уникальный идентификатор. Объект WORLD является сложным, так как содержит объектный атрибут PLAYER, зеркальный группе PlayerWorld объекта PLAYER; остальные его атрибуты простые и однозначные.

    \item \textbf{ROLE} (роль). Для каждой роли необходимо хранить:
    \begin{itemize}
        \item название роли RoleName (1.1);
        \item приоритет роли RolePriority (1.1);
        \item отображаемый префикс в чате ChatPrefix (0.1);
        \item набор предоставляемых прав Rights (1.N);
        \item игроков, которым назначена роль, --- объектный атрибут PLAYER (0.N).
    \end{itemize}
    Атрибут RoleName --- уникальный идентификатор. Роль определяет уровень доступа игрока на сервере; в системе могут быть предусмотрены роли Player, VIP, Moderator, Admin. Атрибут RolePriority предназначен для упорядочивания ролей, обладающих одинаковым набором прав: так, роли Admin и Moderator могут обе предоставлять право наложения и снятия блокировки, однако снять блокировку, наложенную модератором, вправе только администратор, что и определяется значением приоритета соответствующей роли. Объект ROLE является сложным, так как содержит объектный атрибут PLAYER, и одновременно составным, так как содержит многозначный простой атрибут Rights.

    \item \textbf{BAN} (блокировка). Система должна хранить информацию о блокировках игроков:
    \begin{itemize}
        \item заблокированного игрока --- объектный атрибут PLAYER в роли Banned (1.1);
        \item время начала блокировки StartTime (1.1);
        \item причину блокировки Reason (1.1);
        \item дату и время окончания блокировки ExpiresAt (0.1);
        \item признак активности блокировки IsActive (1.1);
        \item администратора, наложившего блокировку, --- объектный атрибут PLAYER в роли Admin (1.1).
    \end{itemize}
    Атрибуты PLAYER (роль Banned) и StartTime являются неуникальными идентификаторами: по игроку и по времени начала можно найти блокировки, но каждый из этих атрибутов по отдельности блокировку однозначно не определяет, а определяет её только их сочетание. Объект BAN является ассоциативным: он сопоставляет двух игроков --- заблокированного и администратора --- и содержит дополнительную информацию об этом сопоставлении (время начала, причину, срок действия и признак активности).
\end{enumerate}


\begin{figure}[ht]
    \centering
    \includegraphics[width=0.87\textwidth]{/home/devone/DatabaseLabs/lab2/DB_Lab2.drawio.png}
    \caption{Семантическая модель объектов}
    \label{fig:som-model}
\end{figure}

\subsection{Обоснование выбора кардинальных связей}

Связи между объектами модели, представленной на \hyperref[fig:som-model]{рисунке~\ref*{fig:som-model}}, задаются парами зеркальных объектных атрибутов, а кардинальное число каждого атрибута записывается в виде <<мин.макс>>. Ниже приведён перечень таких пар с обоснованием выбора кардинальных чисел каждой из них. Тип связи определяется максимальными кардинальными числами обоих атрибутов пары.

\begin{enumerate}
    \item \textbf{ROLE --- PLAYER.} Объектный атрибут ROLE объекта PLAYER имеет кардинальное число 1.1, а зеркальный ему атрибут PLAYER объекта ROLE --- 0.N (связь 1:N). Каждому экземпляру объекта PLAYER соответствует ровно один экземпляр объекта ROLE, поскольку роль определяет единственный уровень доступа игрока на сервере. Одному экземпляру объекта ROLE может соответствовать от нуля до произвольного числа экземпляров объекта PLAYER, так как роль не закреплена за конкретным игроком и может быть либо ещё не назначена ни одному игроку, либо назначена одновременно многим игрокам.

    \item \textbf{PLAYER --- WORLD.} Связь <<многие-ко-многим>> (M:N) реализована групповым атрибутом PlayerWorld объекта PLAYER, внутри которого находится объектный атрибут WORLD (1.1), и зеркальным ему атрибутом PLAYER объекта WORLD (0.N). Каждый экземпляр группы PlayerWorld описывает статистику посещения мира конкретным игроком, поэтому ему соответствует ровно один экземпляр объекта WORLD, а отметка идентификатора у атрибута WORLD не допускает повторения одного и того же мира в статистике игрока. Одному экземпляру объекта PLAYER может соответствовать от нуля (игрок, ни разу не заходивший ни в один мир) до произвольного числа экземпляров группы PlayerWorld (0.N), так как один игрок может посещать произвольное количество миров. Одному экземпляру объекта WORLD, в свою очередь, может соответствовать от нуля (мир, который ещё не посещался ни одним игроком) до произвольного числа экземпляров объекта PLAYER --- по числу посетивших его игроков.

    \item \textbf{PLAYER --- BAN (заблокированный игрок).} Объектный атрибут PLAYER в роли Banned объекта BAN имеет кардинальное число 1.1, а зеркальный ему атрибут BAN в роли Received объекта PLAYER --- 0.N (связь 1:N). Каждый экземпляр объекта BAN соответствует ровно одному экземпляру объекта PLAYER, так как содержит уникальные для данной блокировки сведения --- причину и срок действия. Экземпляру объекта PLAYER может соответствовать от нуля (игрок без нарушений) до произвольного числа экземпляров объекта BAN, поскольку система регистрирует как действующие, так и завершённые блокировки игрока и хранит их историю. Одновременное наложение на игрока нескольких действующих блокировок системой не предусмотрено, что обеспечивается признаком IsActive, а не кардинальным числом.

    \item \textbf{PLAYER --- BAN (администратор, наложивший блокировку).} Объектный атрибут PLAYER в роли Admin объекта BAN имеет кардинальное число 1.1, а зеркальный ему атрибут BAN в роли Issued объекта PLAYER --- 0.N (связь 1:N). Каждому экземпляру объекта BAN соответствует ровно один экземпляр объекта PLAYER, выступающий в роли администратора, наложившего блокировку, так как коллективное наложение блокировки системой не предусмотрено. Одному экземпляру объекта PLAYER, обладающему соответствующими правами доступа, может соответствовать от нуля до произвольного числа наложенных им экземпляров объекта BAN.
\end{enumerate}

Кардинальные числа остальных атрибутов выбраны следующим образом:
\begin{itemize}
    \item идентификаторы и простые атрибуты, не упомянутые ниже, имеют кардинальное число 1.1: их значение известно в момент появления экземпляра объекта и является единственным (например, игрок появляется в системе при первом входе на сервер, поэтому FirstSeenAt и LastSeenAt известны сразу);
    \item групповой атрибут LastPosition и входящие в него атрибуты X, Y и Z имеют кардинальное число 1.1: позиция задана ровно тремя координатами и известна с первого посещения мира;
    \item атрибут ChatPrefix имеет кардинальное число 0.1: не каждая роль отображается в чате с префиксом, но префикс не может быть больше одного;
    \item атрибут Rights имеет кардинальное число 1.N: роль определяет уровень доступа, поэтому предоставляет хотя бы одно право, а число прав в наборе не ограничено;
    \item атрибут ExpiresAt имеет кардинальное число 0.1: блокировка может быть бессрочной и тогда не имеет даты окончания, но срок действия у неё не может быть больше одного;
    \item атрибут IsActive имеет кардинальное число 1.1: блокировка в каждый момент либо действует, либо завершена, а значение признака нельзя вывести из ExpiresAt, так как блокировка может быть снята досрочно.
\end{itemize}

\end{document}
